\subsection{整个薄层区域的解析下界}\label{sec:sliver}
\begin{lemma}[连续相位速度给出的谱隙下界]\label{lem:r8-phase-speed}
对任意 $R\ge1$, $0<u<1/2$, 令 $\varepsilon=R^{-1/2}$,
$\ell=1/2-u$, $w=u/\varepsilon$. 则
\begin{equation}
 G(R,u)\ge\frac{\pi^2}{2\varepsilon(w+\ell)(w+\varepsilon\ell)}.
 \label{eq:r8-global-gap-bound}
\end{equation}
本结论不要求第二个重层相位大于 $\pi/2$.
\end{lemma}
\begin{proof}
对 $k=\sqrt\lambda>0$, 取 $y(0)=0$, $y'(0)=1$ 的射击解.
接口向量为 $(y'(u),ky(u))=(\cos(wk),\varepsilon\sin(wk))$.
经过轻层后, 它旋转了角 $\ell k$, 两分量即 \eqref{eq:polefree-even}--\eqref{eq:polefree-odd}.
令 $A_\varepsilon$ 是向量 $\cos\theta+i\varepsilon\sin\theta$ 的连续展开辐角,
$A_\varepsilon(0)=0$. 向量从不为零, 且
\[
 A_\varepsilon'(\theta)=\frac{\varepsilon}{\cos^2\theta+\varepsilon^2\sin^2\theta}>0,
 \quad A_\varepsilon(\pi/2)=\pi/2,\quad A_\varepsilon(\pi)=\pi.
\]
半区间末端相位为
\[
 \Phi(k)=\ell k+A_\varepsilon(wk),\qquad
 0<\Phi'(k)\le\ell+w/\varepsilon.
\]
它从零严格增加到无穷. 第一 Neumann 根和第一 Dirichlet 根分别满足
$\Phi(k_1)=\pi/2$, $\Phi(k_2)=\pi$.
空间相位在两层中均严格递增, 故这两个半区间解内部为正.
前者偶延拓没有内部零点, 后者奇延拓恰有中点这一个零点;
由振荡定理, $k_1^2,k_2^2$ 确为全区间的前两个特征值.

对这两个实际根积分相位导数, 得
\[
 k_2-k_1\ge\frac{\pi}{2(\ell+w/\varepsilon)}.
\]
又 $\Phi(\pi/(2w))>\pi/2$, 故 $wk_1<\pi/2$.
在第一分支, 因 $0<\varepsilon\le1$,
$A_\varepsilon(\theta)=\arctan(\varepsilon\tan\theta)\le\theta$,
所以 $k_1\ge\pi/[2(\ell+w)]$.
相乘即得
\[
 G=\varepsilon^{-2}(k_2-k_1)(k_2+k_1)
 \ge\frac{\pi^2}{2\varepsilon(w+\ell)(w+\varepsilon\ell)}.
\]
以上分母均正, 没有使用任何全域奇相位分支假设.
\end{proof}

\begin{lemma}[薄层区域, 替代旧网格证书]\label{lem:sliver}
对每个 $R\ge1500$ 及 $0<u\le2/\sqrt R$, 有
\begin{equation}
 G(R,u)>\frac{15\pi^2}{4}>25,\qquad G(R,u)>3\pi^2.
 \label{eq:r8-sliver-margin}
\end{equation}
\end{lemma}
\begin{proof}
此时 $w\le2$, $\ell<1/2$; 由 $38^2<1500$ 得 $\varepsilon<1/38$.
故 \eqref{eq:r8-global-gap-bound} 给出
\[
 G>\frac{\pi^2}{2(1/38)(5/2)(2+1/76)}
 =\frac{2888}{765}\pi^2>\frac{15}{4}\pi^2.
\]
最后一项有理比较等价于 $11552>11475$.
再用 $\pi>3$, 得 $15\pi^2/4>135/4>25$; 另有 $15/4>3$.
这覆盖整个连续区域, 包括任意小的正 $w$, $w=1/2$, 两条旧曲边界及无穷大的 $R$ 尾部.
\end{proof}

